\section{Asymptotes}
\label{sec:asymptotes}

An \textbf{asymptote} is a line that describes how a graph behaves near
some input or far out to the left or right. Limits tell us exactly what
this means. We will focus on horizontal and vertical lines here; see
Section~\ref{sec:slanted-asymptotes} in the appendix for slanted asymptotes
and other ways to describe a function's long-term behavior.

\subsection{Horizontal Asymptotes}

The line $y=L$ is a \textbf{horizontal asymptote} of $f$ if
$$\lim_{x\to\infty}f(x)=L\qquad\text{or}\qquad
\lim_{x\to-\infty}f(x)=L,$$
where $L$ is a finite number. So, to find horizontal asymptotes, compute
these two limits separately. The left and right ends of the graph might
approach different heights, or one end might have no horizontal asymptote.

For example:
\begin{itemize}
\item Remember that $\sqrt{x^2}=|x|$. For $x\ne0$,
$$\frac{x}{\sqrt{x^2+1}}=
\frac{x/|x|}{\sqrt{1+1/x^2}}.$$
The numerator on the right is 1 for $x>0$ and $-1$ for $x<0$, so
$$\lim_{x\to\infty}\frac{x}{\sqrt{x^2+1}}=1,
\qquad
\lim_{x\to-\infty}\frac{x}{\sqrt{x^2+1}}=-1.$$
This graph has two horizontal asymptotes: $y=1$ on the right and $y=-1$
on the left.
\end{itemize}

A horizontal asymptote is not a barrier: the graph can cross it.
The limit only describes what happens far out, not at every input.

\subsection{Vertical Asymptotes}

The line $x=a$ is a \textbf{vertical asymptote} of $f$ if at least one
of its one-sided limits is infinite:
$$\lim_{x\to a^-}f(x)=\pm\infty\qquad\text{or}\qquad
\lim_{x\to a^+}f(x)=\pm\infty.$$
Only one side is needed, and we only approach through inputs in the
domain. The value of $f(a)$, if it is defined, does not affect these limits.

To find vertical asymptotes, look for inputs where the expression becomes
undefined, then check the one-sided limits. For rational functions, factor
and cancel common factors first (away from the excluded inputs), then
check the remaining denominator zeros. A zero in the original denominator
is only a candidate, not a guarantee!

Here are some examples:
\begin{itemize}
\item Let
$$f(x)=\frac{x^2-1}{(x-1)(x-2)}=\frac{x+1}{x-2}
\qquad(x\ne1,2).$$
There are two candidates, but
$$\lim_{x\to1}f(x)=\frac{2}{-1}=-2.$$
So there is a hole at $x=1$, not a vertical asymptote. Near $x=2$, the numerator
$x+1$ is positive, while $x-2$ is negative on the left and positive on
the right. Thus
$$\lim_{x\to2^-}f(x)=-\infty,\qquad
\lim_{x\to2^+}f(x)=\infty,$$
and $x=2$ is a vertical asymptote.

\item Not every vertical asymptote comes from a denominator. For instance,
$$\lim_{x\to0^+}\ln(x)=-\infty,$$
so $x=0$ is a vertical asymptote of $\ln(x)$, even though there is no
left-hand side of its domain near 0.
\end{itemize}
